\section{Phương pháp}
\label{sec:methodology}

\subsection{Ký hiệu và chuẩn hóa sự kiện}
\label{subsec:event-normalization}

Gọi $e=(id,t,u,ip,r,a,s,q)$ là một event đã được parser chấp nhận, lần lượt
gồm định danh, thời gian, người dùng (có thể rỗng), địa chỉ IP, tài nguyên,
hành động, trạng thái và chuỗi request. Parser chỉ nhận CSV hoặc mảng JSON có
đủ các trường của schema; dòng sai định dạng được ghi nhận là lỗi dòng và không
đi vào các bước phát hiện. Trường \texttt{expected\_label} vẫn bắt buộc trong
dataset demo để đánh giá, nhưng không xuất hiện trong hàm phát hiện hay hàm
chấm điểm.

Mỗi phát hiện $f$ mang một loại rủi ro, các event IDs nguồn, evidence và độ tin
cậy $c\in[0,1]$. Bộ chấm điểm chuyển $f$ thành alert $a_f$; do đó alert luôn
giữ liên kết trở lại event thay vì chỉ chứa một nhãn tổng quát.

\subsection{Kiểm tra quyền RBAC}
\label{subsec:rbac-checking}

Với username $u$, tài nguyên $r$ và hành động $a$, predicate quyền của nguyên
mẫu là
\begin{equation}
  P(u,r,a)=1 \Longleftrightarrow \exists \rho:\ (u,\rho)\in UR
  \ \land\ (\rho,(r,a))\in RP\ \land\ \mathrm{active}(u),
  \label{eq:rbac-permission}
\end{equation}
trong đó $UR$ là gán người dùng--vai trò và $RP$ là gán vai trò--quyền trong
SQLite. Predicate này được truy vấn cho các event có \texttt{event\_type} là
\texttt{access} hoặc \texttt{authorization} và có người dùng. Luật
\texttt{RBAC-UNAUTHORIZED} phát hiện khi
\begin{equation}
  \neg P(u,r,a)\ \lor\ (s=\texttt{denied}).
  \label{eq:unauthorized-rule}
\end{equation}
Vì vậy, cả một yêu cầu không được RBAC cấp quyền lẫn một event log ghi nhận bị
từ chối đều được giữ lại làm bằng chứng. Đây là kiểm tra hậu kiểm trong pipeline,
không phải điểm quyết định cấp quyền trực tuyến.

\subsection{Phát hiện theo luật}
\label{subsec:detection-rules}

\textbf{SQL injection.} Đây là luật phát hiện chính thứ nhất. Với mỗi request,
nguyên mẫu kiểm tra tuần tự biểu thức chính quy cho bốn chỉ báo: tautology,
\texttt{union select}, truy vấn xếp chồng và SQL comment. Mỗi event sinh nhiều
nhất một finding, tại mẫu đầu tiên khớp; evidence chứa tên mẫu và phần request
đã cắt ngắn. Các chỉ báo này là heuristic có thể giải thích, không phải bộ phòng
vệ SQL injection hoàn chỉnh \citep{OWASPSQLi}.

\textbf{Dò đoán mật khẩu.} Đây là luật phát hiện chính thứ hai. Các event
\texttt{authentication} được sắp xếp theo thời gian và phân nhóm theo $(u,ip)$.
Trong từng nhóm, một đoạn thất bại kết thúc khi có đăng nhập thành công hoặc khi
khoảng cách giữa hai event liên tiếp lớn hơn cửa sổ $W_p$. Một finding
\texttt{AUTH-PASSWORD-GUESSING} được tạo nếu trong đoạn có một cửa sổ trượt đạt
ít nhất $N_p$ event thất bại. Cấu hình mặc định dùng $N_p=5$ và $W_p=300$ giây.
Finding giữ toàn bộ event IDs của đoạn để evidence có thể được rà soát.

\textbf{Truy cập trái phép.} Luật trong \eqref{eq:unauthorized-rule} là kiểm
tra hỗ trợ: nó làm rõ event \texttt{access}/\texttt{authorization} nào không
được RBAC cấp quyền hoặc đã bị log ghi nhận là từ chối. Các finding từ ba luật
được sắp xếp xác định theo thời gian, loại rủi ro và rule ID trước khi đưa vào
bước chấm điểm.

\subsection{Tín hiệu rủi ro theo ngữ cảnh}
\label{subsec:context-signals}

Khi tùy chọn context-risk được bật, profiler duyệt event theo thời gian và chỉ
học baseline từ các event đã quan sát trước đó. Đây là phần mở rộng để ưu tiên
triage, không làm thay đổi predicate quyền trong \eqref{eq:rbac-permission}.
Bốn tín hiệu có thể giải thích là: (i) \texttt{CTX-AFTER-HOURS} nếu giờ event
nằm ngoài $[h_s,h_e)$; (ii)
\texttt{CTX-NEW-IP} nếu IP chưa từng thấy cho người dùng nhưng baseline IP đã
có; (iii) \texttt{CTX-RARE-RESOURCE} nếu event là access/authorization và tài
nguyên chưa từng được người dùng truy cập thành công; và (iv)
\texttt{CTX-REPEATED-DENIAL} khi đúng event từ chối thứ ba của cùng $(u,ip)$
nằm trong cửa sổ $W_c$. Các giá trị mặc định là $h_s=8$, $h_e=18$ và
$W_c=300$ giây. Sau khi xử lý một event, profiler mới thêm IP vào baseline và
chỉ thêm tài nguyên khi event thành công; nhờ đó event đầu tiên không tự tạo tín
hiệu IP/tài nguyên mới.

\subsection{Chấm điểm và phân mức}
\label{subsec:risk-scoring}

Với finding theo luật, điểm alert được tính xác định bằng
\begin{equation}
  R_f=\min\{100,\ B_k+\operatorname{round}(20c)
  +\min(15,3\max(0,n-1))+15d\},
  \label{eq:rule-score}
\end{equation}
trong đó $B_k\in\{45,40,50\}$ lần lượt cho SQL injection, password guessing
và unauthorized access; $n$ là số lần lặp của finding và $d=1$ khi RBAC từ
chối, ngược lại bằng 0. Với context finding, công thức là
\begin{equation}
  R_c=\min\{100,\ 25+\operatorname{round}(20c)+b_s\},
  \label{eq:context-score}
\end{equation}
trong đó $b_s$ là bonus theo signal (mặc định 10 cho IP mới, 8 cho ngoài giờ và
tài nguyên hiếm, 12 cho từ chối lặp lại). Theo cấu hình, mức Medium, High và
Critical bắt đầu tại 30, 60 và 85; điểm thấp hơn 30 là Low.

\subsection{Gom nhóm incident}
\label{subsec:incident-grouping}

Các alert chỉ được gom thành incident khi context-risk được bật. Như
Hình~\ref{fig:event-alert-incident}, thuật toán phân vùng alert theo $(u,ip)$,
sắp xếp mỗi phân vùng theo timestamp và bắt đầu đoạn mới nếu khoảng cách với
alert ngay trước lớn hơn $W_c$. Với một đoạn $G$, điểm incident là
\begin{equation}
  R_I=\min\{100,\max_{a\in G}R_a + 15\,\mathbf{1}(|G|>1)\}.
  \label{eq:incident-score}
\end{equation}
Incident thu thập theo thứ tự không lặp các event IDs, alert IDs, evidence,
loại rủi ro và signal ngữ cảnh của đoạn. Hai trường \texttt{start\_time} và
\texttt{end\_time} là thời gian alert đầu/cuối của đoạn, không phải một timeline
pháp chứng chi tiết.

\begin{figure}[t]
  \centering
  \resizebox{0.98\linewidth}{!}{% Quy tắc build_incidents trong src/rbac_guard/context.py.
\begin{tikzpicture}[
  font=\footnotesize,
  >=Latex,
  node distance=8mm and 11mm,
  entity/.style={draw, rounded corners=2pt, align=center, minimum height=12mm,
    text width=29mm, fill=blue!7},
  grouping/.style={draw, rounded corners=2pt, align=center, minimum height=13mm,
    text width=34mm, fill=orange!14},
  result/.style={draw, rounded corners=2pt, align=center, minimum height=15mm,
    text width=34mm, fill=green!11},
  flow/.style={->, line width=.55pt}
]
  \node[entity] (events) {Event\\ID, thời gian, user, IP,\\request, trạng thái};
  \node[entity, right=of events] (alerts) {Alert\\rule/signal, evidence,\\điểm, mức độ, event IDs};
  \node[grouping, right=of alerts] (partition) {Phân vùng theo\\$(user, IP)$; sắp xếp\\theo thời gian};
  \node[grouping, below=of partition] (segment) {Tạo đoạn mới khi\\khoảng cách $> W$;\\$W=300$ giây trong cấu hình};
  \node[result, right=16mm of segment] (incident) {Incident\\IDs alert/event, evidence,\\start/end, điểm, mức độ};

  \draw[flow] (events) -- node[above, align=center] {luật và\\ngữ cảnh} (alerts);
  \draw[flow] (alerts) -- (partition);
  \draw[flow] (partition) -- (segment);
  \draw[flow] (segment) -- node[above, align=center, font=\scriptsize] {max alert +\\bonus chuỗi} (incident);
\end{tikzpicture}
}
  \caption{Phần mở rộng context-risk: luồng từ event sang alert và incident.
  Alert được phân vùng theo cặp người dùng--IP rồi sắp xếp theo thời gian. Trong
  mỗi phân vùng, một khoảng cách lớn hơn cửa sổ $W$ tạo incident mới; các trường
  thời gian chỉ là start/end của đoạn, không phải timeline pháp chứng chi tiết.}
  \label{fig:event-alert-incident}
\end{figure}
